\section{Complete physical roof variation without a margin cutoff}
\label{sec:rbv-complete}

The full-source weak roof estimates of Section~\ref{sec:fr-full-roof}
leave open a pointwise concentration at a common physical roof.
We now bound that common-roof density by summing its actual interior
variation and all geometric entrance and exit charges.  This is a
directional variation argument in the physical roof variable.  No
variation in the Gauss angle is asserted.  The local critical
remainders are included in the exact interior calculation; they are
not summed after a word-dependent Taylor approximation.

Set $a=5/6$, and retain the original regular endpoint cap
\[
 \mathcal C_{m,R}=\{x: |p_0(x)|,|p_m(x)|<a,\quad
                     x\hbox{ is an actual regular }m\hbox{-flight record}\}.
\]
Its boundary has collision area zero, so the strict convention has
the same bulk measure as the closed convention in 314.  Every
original exact event lies in the corresponding closed cap.  Define
the positive majorant, on the unchanged unit-speed roof line, by
\begin{equation}\label{eq:rbv-majorant}
 \overline\mu_{m,R}
      =(L_{m,R})_*\bigl(\one_{\mathcal C_{m,R}}\,d\nu/c_*\bigr).
\end{equation}
This is not a new probability or a new conditioning event: its
mass is at most $c_*^{-1}$, and the original initial-section
indicator is restored by positive restriction.  In particular,
for every original exact record and every Borel $0\le H\le1$,
\begin{equation}\label{eq:rbv-domination}
 (L_{m,R})_*(\one_{\Pi_{n,k,m,R}}H\,d\nu_R^*)
                  \le\overline\mu_{m,R}.
\end{equation}
Neither side changes the collision count, physical clock or labels.

Keep the already fixed constants $\lambda_\circ,\Lambda_\circ,
C_\omega,C_F,J_\times,C_{\rm mom},C_{\rm ab}$ and
$\sigma=47/59$.  Define
\begin{align}
 C_{\rm int}
   &=\frac{\Lambda_\circ}{\lambda_\circ}
                 \left(C_\omega+\frac{2C_F}{\lambda_\circ}\right),
                                                    \label{eq:rbv-constants}\\
 D_{\rm roof}
   &=\frac{C_{\rm int}J_\times}{c_*}
       +\frac{2C_{\rm ab}}{1-\sigma}
       +\frac{\Lambda_\circ C_{\rm mom}}{2c_*R_-^2}
       +\frac{U_2 C_{\rm mom}}{2c_*},
 & U_2&=320/9,\qquad C_{\rm roof}=D_{\rm roof}/2. \notag
\end{align}
These constants are fixed by the original physical family and the
native homogeneous curve class, before every collision, radius,
label or roof supremum.  They contain no whole-word pressure or
comparison-capacity assumption.

\subsection{The unweighted interior and the actual cap boundaries}

\begin{lemma}[The full regular interior has finite roof variation]
\label{lem:rbv-interior}
On a regular fixed-Gauss-angle physical segment put
$t=F$, $g=\nabla_qF$, $z=|g|>0$ and
$A=\omega/(4\pi c_*\chi)$.  Then
\begin{equation}\label{eq:rbv-interior-pointwise}
       |\partial_t A|\le C_{\rm int} A I_m/z.
\end{equation}
On the complete cap source, with all actual words summed,
\begin{equation}\label{eq:rbv-interior-total}
 \int_0^{2\pi}\sum_w\int_{\mathcal D_w(\theta)}
                 |\partial_t A_w(t,\theta)|\,dt\,d\theta
                         \le C_{\rm int}J_\times/c_*.
\end{equation}
Here $\mathcal D_w(\theta)$ is the open set of actual cap records
of the indicated word on that Gauss ray.  Missing physical intervals
are not filled in.
\end{lemma}
\begin{proof}
The proof of Theorem~\ref{thm:el-regular}, before differentiating
its guard, gives
\[
 |\partial_t\log A|
 \le\frac{\Lambda_\circ}{\lambda_\circ z}
        \left(C_\omega I_m+\frac{2C_F}{\lambda_\circ}\right).
\]
The endpoint bounds used there depend only on
$|p_0|,|p_m|\le9/10$ and hence hold on this cap.  Since $I_m\ge1$,
this proves \eqref{eq:rbv-interior-pointwise}, with no guard
derivative and no factor $\varepsilon^{-1}$.

The exact endpoint-coordinate coarea formula in
Theorem~\ref{thm:v111-unique} and
Proposition~\ref{prop:v112-exhaustion} applies to the complete
regular cap just as to its original exact-event subsets: the
uniqueness proof uses the physical word, and not the section test.
The critical set has zero collision area.  Consequently
\[
 \int_0^{2\pi}\sum_w\int_{\mathcal D_w(\theta)}
                    A_w I_m/z\,dt\,d\theta
    =c_*^{-1}\int_{\mathcal C_{m,R}}Z_m\,d\nu
    \le J_\times/c_*.
\]
The last inequality is the complete physical multiplier theorem
\ref{thm:sm-full-multiplier}; no independent product law is used.
All equalities follow first on compact regular subsegments, then
by nonnegative exhaustion.  This proves the asserted full sum.
\end{proof}

\begin{lemma}[The four endpoint-cap charges]
\label{lem:rbv-cap}
The complete positive roof--angle entrance/exit charge on the four
regular sides $p_0=\pm a$, $p_m=\pm a$, with the other endpoint
momentum in $[-a,a]$, is at most
\begin{equation}\label{eq:rbv-cap}
                B_{\rm cap}
                    =\frac{\Lambda_\circ C_{\rm mom}}
                                  {2c_*R_-^2}.
\end{equation}
All actual physical words are included before any original record
restriction.  Visited singular intersections have the separate
zero trace stated below.
\end{lemma}
\begin{proof}
Consider $p_0=a$ and write $v=p_m$.  Since $a>0$, the Gauss
angle of $(-a,v)$ has absolute differential
$a(a^2+v^2)^{-1}dv$.  The exact optical identity of
Corollary~\ref{cor:bf-wronskian} is
\[
 A=\frac{u_x}{4\pi c_*h_x},\qquad
 h_x=R^2e^{\mathsf T}G^{-1}e,\qquad
 A\le\frac{\Lambda_\circ}{4\pi c_*R_-^2}u_x.
\]
The complete physical momentum-root bound of
Theorem~\ref{thm:mc-kernel} applies at the common pair $(a,v)$.
Thus this side has charge at most
\[
 \frac{\Lambda_\circ C_{\rm mom}}{4\pi c_*R_-^2}
                  \int_{-a}^a\frac{a}{a^2+v^2}\,dv
          =\frac{\Lambda_\circ C_{\rm mom}}{8c_*R_-^2}.
\]
The remaining sides have exactly the same integral, either by
changing the sign or interchanging the two endpoint coordinates.
Summing four positive bounds proves \eqref{eq:rbv-cap}.
Cap corners are assigned once; they cannot add a charge by
duplicating a coordinate chart.  This argument evaluates the
proved pointwise root sum on its actual regular roots.  It does
not restrict an arbitrary version of a two-dimensional density
to a null line.
\end{proof}

\subsection{Exhausting all physical entrances and exits}

\begin{lemma}[Complete fixed-count physical endpoint accounting]
\label{lem:rbv-endpoints}
Fix $R,m$.  Up to a null set of Gauss angles, every endpoint of
a maximal regular fixed-angle interval in the cap has one of the
following actual sources: a cap side; an unvisited-disk tangency;
an interior visited grazing collision; or $g=0$.  The complete
positive charge of the first two types is bounded by
\begin{equation}\label{eq:rbv-boundary-total}
                   B_{\rm cap}+\frac{2C_{\rm ab}}{1-\sigma}.
\end{equation}
The third type has zero coarea endpoint trace.  The complete
incoming charge at $g=0$ is at most
\begin{equation}\label{eq:rbv-critical-total}
                         B_{\rm crit}
                             =\frac{U_2C_{\rm mom}}{2c_*}.
\end{equation}
There is no additional charge at a homogeneity cut, angular seam,
nearest-obstacle switch or auxiliary coordinate boundary.
\end{lemma}
\begin{proof}
At fixed $R,m$ the lifted word family is finite, by the actual
finite horizon and lattice bound.  The global uniqueness theorem
\ref{thm:v111-unique} makes the regular roof--angle image for one
word an open set, even when its physical domain is disconnected.
At a fixed angle its maximal components are disjoint open roof
intervals.  The endpoint-limit argument in
Theorem~\ref{thm:v109-traces} applies without a section selector:
both endpoint angles, every designated contact, flight, incidence
and third-disk clearance have limits.  Flight lengths stay at
least $\tau_-$.  The cap keeps both endpoint incidences positive.
If the limit has $g\ne0$, no visited grazing, positive third-disk
clearance and a strict cap inequality, the local physical map and
roof--angle inverse extend through it.  It therefore is not an
endpoint of a maximal component.  This proves the list.  For a
disjoint charge assignment, assign visited-grazing limits first,
then the remaining critical limits.  Among noncritical limits,
assign an active unvisited constraint before a cap constraint.
Thus a residual cap arc consists of actual regular roots of the
momentum map, as required in Lemma~\ref{lem:rbv-cap}.

Here is a precise way to account for simultaneous boundaries and
coordinate limits.  The closure of the finite-word configurations
in the cap is compact.  The formal analytic endpoint charts of
Theorem~\ref{thm:pt-analytic}, including their visited-grazing
limits, give a finite local cover for this fixed $R,m$.  On each
such chart, actual admissibility is a finite Boolean combination
of analytic sign conditions: the signed visited incidences and
the distances and foot tests for finitely many candidate disks.
Endpoint-distance cases are retained as well.  This finite list
follows from the horizon, not a sampled list of encountered disks.
The cap adds four analytic inequalities.

Where $g\ne0$, the continuous positive reduced Hessian in 260
makes $(t,\theta)$ local analytic coordinates, also for the formal
calculation at a one-sided limit.  On a fixed-angle slice a
nonconstant analytic boundary function has isolated zeros on
every compact subinterval.  If it is identically zero on that
slice it imposes a constant truth value there; exceptional angular
levels of this type may be discarded when the function is not
identically zero in a neighborhood.  Indeed an accumulating
sequence of such angular levels would make the two-variable
analytic function identically zero.  Identical constraints are
removed from the local Boolean description.  Thus Fubini reduces
the calculation to ordinary one-dimensional intervals.  A smooth
partition of unity on the finite chart cover introduces artificial
terms whose sum is zero before absolute values are taken.  Equivalently
one may use the global one-copy word image and cancel its artificial
subinterval endpoints first.  No new physical component is charged.

At a nongrazing unvisited boundary, its closest foot is strictly
inside the indicated flight: a third disk is separated from either
incident contact by the positive disk gap.  The marked-foot
standoff therefore applies.  The signed line-clearance differential
is nonzero by the actual determinant in
Theorem~\ref{thm:bf-foot-determinant}.  A regular boundary arc
contributes its trace times $|d\theta|$, exactly as in
Theorem~\ref{thm:bf-flux}.  An arc on which the angle derivative
vanishes has zero angular image, by the one-dimensional argument
in that theorem.  Intersections or endpoints with zero angular
image introduce no extra interval-endpoint integral.  If several
disk constraints are active, assign the physical boundary germ to
the first of them in a fixed ordering; the full positive marked
sum is an upper bound.  Coincident boundary arcs are assigned once,
and isolated corners are treated by the same fixed-angle slicing.
All actual marks, words and both impact signs are already present
in Theorem~\ref{thm:ab-all-offset} at $u_0=0$.

The proof of that theorem starts from all regular angular records
and uses only the endpoint cap $|p_0|,|p_m|\le9/10$; its later
exact-record tests are optional positive selectors.  It is therefore
valid for the present smaller cap, without imposing an initial
section.  Also its proof bounds $A$ itself: on this zero-clearance
boundary the original $W^\varepsilon$ equals one.  Summing its
finite-count estimate gives
\[
 \sum_{j=0}^{m-1}\mathfrak B^{\rm reg}_{j,0}
    \le C_{\rm ab}\sum_{j=0}^{m-1}
                 (\sigma^j+\sigma^{m-1-j})
    \le\frac{2C_{\rm ab}}{1-\sigma}.
\]
Together with Lemma~\ref{lem:rbv-cap}, this proves
\eqref{eq:rbv-boundary-total}.

At any actually visited interior grazing limit the marked optical
bound in 260 gives $\omega\to0$, whereas
$\chi\ge\lambda_\circ$.  Hence $A\to0$, even if $g\to0$ at
the same endpoint.  This is the inherited zero \emph{trace}
proof, not an appeal to collision-area nullness.  It is applied
before an integral over endpoint angles.  At fixed count the
critical limiting configuration of a word, if it exists, is
unique.  The support inequalities of
Lemma~\ref{lem:v111-support} pass to the limit; normal endpoints
force identical endpoint contacts, and equality in the length
support inequalities then forces the same successive entry or
tangent contacts.  Thus there are only finitely many such critical
configurations at this stage.

For a critical germ with positive visited incidences, the local
physical side is analytic.  The cap and all visited-incidence
tests are strict there.  Every active physical boundary is
therefore an unvisited clearance with the nonzero differential
already proved above, so the specific nonzero-differential
hypothesis of Theorem~\ref{thm:mc-critical-sectors} is satisfied.
In the gradient coordinates, the first
nonzero homogeneous Taylor term of each active nonconstant
inequality has only finitely many zero rays.  Off those rays its
sign is constant along every sufficiently short ray.  This proves
the angular sector limit, including simultaneous tangencies and
collapsed slices, by exactly the analytic argument in
Lemma~\ref{pc:representatives}.  There is no section or arbitrary
Borel selector in the majorant being differentiated.  Its incoming
angular trace is exactly $\mathfrak J(x)\beta_y(x)$, with the
normalized angular change of variables of
Lemma~\ref{lem:mc-angle}.  The complete one-copy sector bound of
Theorem~\ref{thm:mc-critical-sectors} gives
\eqref{eq:rbv-critical-total}; it already includes regular critical
points.  At a visited-grazing critical germ the incoming trace is
zero by the preceding optical argument.  One may verify this by
excising $|g|<\delta$ on its fixed local chart: $A\to0$ uniformly
there, while the angular range is at most $2\pi$.  A fixed finite
family permits the local limit first, without a collision-uniform
chart radius.

The positive Hessian makes every critical endpoint incoming on its
fixed ray; no second outgoing critical charge is added.  All local
limits in this proof occur at fixed $R,m$, after which the displayed
uniform bounds are applied.  The complete regular interior and
regular boundary are kept in their exact measures, so no uniform
Taylor remainder or exchange of long-time and local limits is used.
\end{proof}

\begin{theorem}[Uniform full-source variation and original-roof height]
\label{thm:rbv-height}
For every $m\ge1$ and $R\in[9/20,47/100]$, the measure
\eqref{eq:rbv-majorant} has a nonnegative density
$\overline p_{m,R}\in BV(\mathbb R)$ with
\begin{equation}\label{eq:rbv-height}
 \|D\overline p_{m,R}\|_{\rm TV}\le D_{\rm roof},\qquad
 \|\overline p_{m,R}\|_\infty\le C_{\rm roof},\qquad
 \|\overline p_{m,R}\|_1\le c_*^{-1}.
\end{equation}
In particular, for every original exact return, displacement and
collision label, every Borel $0\le H\le1$, and every Borel physical
roof set $S$, the unchanged source satisfies
\begin{equation}\label{eq:rbv-exact-source}
 \nu_R^*\bigl(\one_{\Pi_{n,k,m,R}}H\one_{\{L_{m,R}\in S\}}\bigr)
                         \le C_{\rm roof}|S|.
\end{equation}
The corresponding density is at most $C_{\rm roof}$ almost
everywhere.  The complete physical sum precedes
$\sup_{R,m,n,k}$ and the roof essential supremum.
\end{theorem}
\begin{proof}
The exact coarea formula gives the density
\[
 \overline p(t)=\int_0^{2\pi}\sum_w
               \one_{\mathcal D_w}(t,\theta)A_w(t,\theta)\,d\theta.
\]
It is integrable with mass at most $c_*^{-1}$.  On each maximal
regular interval integrate $A_w\psi'$ by parts, for
$\psi\in C_c^1(\mathbb R)$.  The endpoint trace and absolute
continuity follow from the proof in 260 and 289 with weight one.
First perform the calculation on a finite interval exhaustion.
The absolute sums of interior derivatives and endpoint charges are
finite after angular integration by the two preceding lemmas.
The integral of $A_w|\psi'|$ is finite by the source mass.
Consequently the exhaustion passes through the integrated identity
and gives
\[
 \left|\int\overline p(t)\psi'(t)\,dt\right|
                  \le D_{\rm roof}\|\psi\|_\infty.
\]
More explicitly, the derivative is the signed pushforward of the
integrable interior derivatives and the finite entrance and exit
trace measures just constructed.  Thus it is a finite signed
measure, and the inequality is the first bound in
\eqref{eq:rbv-height}.

The density is supported in $[m\tau_-,m\tau_+]$.  Its standard
one-dimensional bounded-variation representative is zero outside
that interval.  To reach any nonnegative value $v$ from zero and
return to zero requires total variation at least $2v$.
Approximation by Lebesgue points therefore gives
$2\|\overline p\|_\infty\le\|D\overline p\|_{\rm TV}$.
This proves the remaining bounds.  Finally
\eqref{eq:rbv-domination} and positivity give
\eqref{eq:rbv-exact-source} and its density version.  At each fixed
radius the null sets can be chosen for the countable original
labels; no common null set across all real radii is asserted.
\end{proof}

\begin{corollary}[The unchanged bad guard and exact roof Fourier payment]
\label{cor:rbv-guard}
For the full cap majorant weighted by the original
$W^\varepsilon=1-H^\varepsilon$, $\varepsilon>0$, the same result
holds with
\begin{equation}\label{eq:rbv-guard-constant}
 D_{\rm roof}^\varepsilon
  =\frac{C_{\rm reg}(1+\varepsilon^{-1})J_\times}{c_*}
    +\frac{2C_{\rm ab}}{1-\sigma}+B_{\rm cap}+B_{\rm crit}.
\end{equation}
Every original exact-event submeasure of this guarded majorant has
roof density at most $D_{\rm roof}^\varepsilon/2$.
For either of the two full cap majorants, if its variation bound
is denoted by $D$ and its density by $p$, then
\begin{equation}\label{eq:rbv-fourier}
 \left|\int e^{ist}p(t)\,dt\right|
            \le\min\{c_*^{-1},D/|s|\}\quad(s\ne0),\qquad
 \|p(\cdot+h)-p\|_1\le D|h|.
\end{equation}
The Fourier assertion in this corollary is for the specified full
majorants.  An arbitrary measurable restriction does not inherit
their variation or Fourier bound from positive domination alone.
\end{corollary}
\begin{proof}
Use the exact regular-variation inequality of
Theorem~\ref{thm:el-regular} in the interior.  The unchanged guard
is continuous on each one-sided finite-word closure and between
zero and one.  Every boundary charge is therefore bounded by the
same unweighted charge; visited-grazing traces remain zero.
The preceding integration by parts proves
\eqref{eq:rbv-guard-constant}.  There is no change in the guard
depth schedule or its physical collision and clearance factors.

For the Fourier bound, integrate by parts against the finite measure
$Dp$, using compact support.  For $s\ne0$ the resulting factor is
$1/(is)$; its absolute value gives $D/|s|$.  The mass gives the
other bound.  For translations, the bounded-variation representative
satisfies $|p(t+h)-p(t)|\le |Dp|$ on the interval between $t$ and
$t+h$, up to endpoint conventions on null sets.  Tonelli integration
over $t$ gives $D|h|$.
\end{proof}

\paragraph{The exact gain and the remaining normalized question.}
Theorem~\ref{thm:rbv-height} is an unscaled, collision-uniform
essential-height estimate on the complete physical source.  It
includes arbitrarily small positive incidence and clearance,
all relevant words, the branch-dependent radii selected by one
common physical roof, and the exact nonleading coarea contributions.
It strengthens the full-source weak norm in 314 by a different
geometric mechanism.  The inherited results 260--322 remain intact.

It does not give an $o(m^{-2})$ height for a small selected source:
a bounded density of small mass can still have a narrow positive
spike.  Nor does differentiation of the cap majorant differentiate
the original section or a pairing selector, establish abundance of
the same-roof comparison pool, or provide a uniform further
conditioning denominator.  At the original application width
$\varepsilon=\varepsilon(B)$, the bound in
\eqref{eq:rbv-guard-constant} is fixed before the collision limit.
No exchange of that limit with $B\to\infty$ is made.  The exact
shortage in \eqref{eq:sc-original-obligation}, the joint twisted
spectral interfaces, and the full raw local limit and four-coordinate
inverse theorem retain their original unresolved status.
